%! Author = chtompki
%! Date = 1/14/26
\documentclass[12pt]{amsart}
% Default margins are too wide all the way around. I reset them here
\setlength{\hoffset}{-.75in}
\setlength{\textwidth}{6.5in}
\setlength{\voffset}{-.5in}
\setlength{\textheight}{9.0in}
\usepackage{amsmath}
\usepackage{leftindex}
\usepackage{amssymb}
\usepackage{amsthm}
\usepackage{enumitem}
\usepackage{setspace}

\newenvironment{solution}
{\begin{proof}[Solution]}
{\end{proof}}

\title{Combinatorics\\Assignment 4}
\author{Rob Tompkins}
\date{\today}

\begin{document}
    \maketitle
    \date{\today}
    \begin{onehalfspace}
        1. How many paths of length $k$, $P_k$, are there in the complete graph on $n$
        edges, $K_n$?

        \begin{solution}
            Let $V=\{v_0,\ldots,v_{n-1}\}$ to bet the vertices of a graph.
            Any ordered $k$-tuple chosen from $V$, like $(\underbrace{v_0,\ldots,v_8,\ldots,v_k, \ldots, v_n}_{k\text{-tuples}})$ represents
            a path.
            For a specific set of indices $i_0,\ldots,i_{k-1}$ we can represent an arbitrary path of length $k$ as
            $(v_{i_0},v_{i_1},\ldots,v_{i_{k-1}})$. Now consider the path $(v_{i_{k-1}}, \ldots, v_{i_1}, v_{i_0})$
            which is the same path as before but from the other direction.
            Thus, we see that we list every path twice when trying to count, once forwards and once backwards.

            So we begin by choosing $k$ vertices from our original $n$ vertices in $\binom{n}{k}$ ways.
            Then, we choose our vertex ordering, which is a permutation in our index set, chosen in $k!$ ways,
            but like we said before we'll be double counting, so it turns out that we need $\frac{k!}{2}$.
            Thus, the total count follows as
            \begin{displaymath}
                \# = \binom{n}{k}\frac{k!}{2} = \frac{n!}{2(n-k)!}
            \end{displaymath}
        \end{solution}

        2. How many 1-factors are there in $K_{2n}?$

        \begin{solution}
            There are $(2n)!$ ways to arrange our $2n$ vertices.
            Let $V=\{v_1,\ldots , v_{2n}\}$ be our vertex set, and let
            the indices $i_1,\ldots, i_{2n}$ represent an arrangement
            $(v_{i_1},v_{i_2}, \ldots,v_{i_{2n}})$.

            Now let's pair off the vertices in our arrangement in the following fashion
            $((v_{i_1},v_{i_2}),(v_{i_3},v_{i_4}),\ldots,(v_{2n-1},v_{2n}))$.
            Now we count this construction. There are $(2n)!$ ways to choose the indices on
            our vertices, but in our pairing as we've listed, the order of the $n$-pairs does not matter
            in the slightest so we will need to divide by $n!$ in the end.
            The ordering of the pair of vertices does not matter which happens in $2^n$ ways. Thus
            our count comes out as
            \begin{displaymath}
                \# = \frac{(2n)!}{2^n(n)!}.
            \end{displaymath}

        \end{solution}
    \end{onehalfspace}
\end{document}
